% titel: Blatt 2, A4: Innenwinkelsumme im n-Eck

% Teil I: Werkzeuge, kein Lösungsweg.

Die Summe der Innenwinkel eines $n$-Ecks (Eckpunkte auf einer Kreislinie,
$x\ge3$) per vollständiger Induktion als $(n-2)\cdot180^\circ$ zeigen.

\begin{werkzeugkasten}[Vollständige Induktion]
Beweisschema: Induktionsanfang ($n=3$: Dreieck $180^\circ$) und
Induktionsschritt ($n\to n+1$: Winkel eines $(n+1)$-Ecks aus dem eines
$n$-Ecks ableiten).
\end{werkzeugkasten}

\begin{werkzeugkasten}[Triangulierung]
Ein konvexes $(n+1)$-Eck zerlegt man durch eine Diagonale in ein Dreieck
und ein $n$-Eck. Die Innenwinkel des großen Polygons setzen sich aus der
Summe des $n$-Ecks plus der $180^\circ$ des Dreiecks zusammen:
\begin{equation}
  S(n+1)=S(n)+180^\circ .
\end{equation}
Da die neuen Punkte auf einer Kreislinie liegen, sind die Polygone konvex
-- jede Diagonale liegt im Inneren.
\end{werkzeugkasten}

\begin{achtung}[Darauf achten]
\begin{itemize}
  \item Die Diagonale muss das Polygon in ein Dreieck und ein $n$-Eck
        zerlegen, \emph{nicht} in zwei beliebige Teile.
  \item Induktionsanfang bei $n=3$: Dreieckssumme $180^\circ=(3-2)\cdot180^\circ$.
  \item Im Schritt die konkrete Summe $(n-2)180^\circ+180^\circ=(n-1)180^\circ
        =((n+1)-2)180^\circ$ zeigen, nicht nur ``$+180^\circ$'' schreiben.
\end{itemize}
\end{achtung}
